\begin{animation}[id="graph-mst", label="Minimum Spanning Tree (Kruskal)"]
\shape{G}{Graph}{nodes=["A","B","C","D","E"], edges=[("A","B"),("A","C"),("B","C"),("B","D"),("C","D"),("C","E"),("D","E")]}

\step
\narrate{Graph with 5 nodes. We will build MST by adding cheapest edges that do not form a cycle.}

\step
\recolor{G.edge[(A,B)]}{state=good}
\recolor{G.node[A]}{state=done}
\recolor{G.node[B]}{state=done}
\narrate{Edge A-B (weight 1) is cheapest. Add to MST. Components: (A,B), (C), (D), (E).}

\step
\recolor{G.edge[(C,E)]}{state=good}
\recolor{G.node[C]}{state=done}
\recolor{G.node[E]}{state=done}
\narrate{Edge C-E (weight 2) is next cheapest. Add to MST. Components: (A,B), (C,E), (D).}

\step
\recolor{G.edge[(B,D)]}{state=good}
\recolor{G.node[D]}{state=done}
\narrate{Edge B-D (weight 3). Add to MST. Components: (A,B,D), (C,E).}

\step
\recolor{G.edge[(A,C)]}{state=good}
\recolor{G.edge[(B,C)]}{state=dim}
\recolor{G.edge[(C,D)]}{state=dim}
\recolor{G.edge[(D,E)]}{state=dim}
\narrate{Edge A-C (weight 4) connects the two components. MST complete with 4 edges. Remaining edges skipped.}
\end{animation}
